\documentclass[10pt,a4paper]{amsart}
\usepackage{amsmath,amssymb,amsthm,graphicx}
% Retained minimal preamble for the Stacks Project chapter sites-cohomology.tex (Cohomology on Sites).
% Verbatim excerpt of the chapter preamble: Xy-pic setup (native preamble lines 14-18) and the native
% theorem declarations with the mathematical macros (native preamble lines 154-217). Copied mechanically;
% nothing added, removed or redefined. The wrapper class and the amsmath/amssymb/amsthm/graphicx base
% packages are supplied by the assembler, so the class/externaldocument part of the native preamble is omitted.
% For commutative diagrams we use Xy-pic
\usepackage[all]{xy}

% We use 2cell for 2-commutative diagrams.
\xyoption{2cell}
% Theorem environments.
%
\theoremstyle{plain}
\newtheorem{theorem}[subsection]{Theorem}
\newtheorem{proposition}[subsection]{Proposition}
\newtheorem{lemma}[subsection]{Lemma}

\theoremstyle{definition}
\newtheorem{definition}[subsection]{Definition}
\newtheorem{example}[subsection]{Example}
\newtheorem{exercise}[subsection]{Exercise}
\newtheorem{situation}[subsection]{Situation}

\theoremstyle{remark}
\newtheorem{remark}[subsection]{Remark}
\newtheorem{remarks}[subsection]{Remarks}

\numberwithin{equation}{subsection}

% Macros
%
\def\lim{\mathop{\mathrm{lim}}\nolimits}
\def\colim{\mathop{\mathrm{colim}}\nolimits}
\def\Spec{\mathop{\mathrm{Spec}}}
\def\Hom{\mathop{\mathrm{Hom}}\nolimits}
\def\Ext{\mathop{\mathrm{Ext}}\nolimits}
\def\SheafHom{\mathop{\mathcal{H}\!\mathit{om}}\nolimits}
\def\SheafExt{\mathop{\mathcal{E}\!\mathit{xt}}\nolimits}
\def\Sch{\mathit{Sch}}
\def\Mor{\mathop{\mathrm{Mor}}\nolimits}
\def\Ob{\mathop{\mathrm{Ob}}\nolimits}
\def\Sh{\mathop{\mathit{Sh}}\nolimits}
\def\NL{\mathop{N\!L}\nolimits}
\def\CH{\mathop{\mathrm{CH}}\nolimits}
\def\proetale{{pro\text{-}\acute{e}tale}}
\def\etale{{\acute{e}tale}}
\def\QCoh{\mathit{QCoh}}
\def\Ker{\mathop{\mathrm{Ker}}}
\def\Im{\mathop{\mathrm{Im}}}
\def\Coker{\mathop{\mathrm{Coker}}}
\def\Coim{\mathop{\mathrm{Coim}}}

% Boxtimes
%
\DeclareMathSymbol{\boxtimes}{\mathbin}{AMSa}{"02}

%
% Macros for moduli stacks/spaces
%
\def\QCohstack{\mathcal{QC}\!\mathit{oh}}
\def\Cohstack{\mathcal{C}\!\mathit{oh}}
\def\Spacesstack{\mathcal{S}\!\mathit{paces}}
\def\Quotfunctor{\mathrm{Quot}}
\def\Hilbfunctor{\mathrm{Hilb}}
\def\Curvesstack{\mathcal{C}\!\mathit{urves}}
\def\Polarizedstack{\mathcal{P}\!\mathit{olarized}}
\def\Complexesstack{\mathcal{C}\!\mathit{omplexes}}
% \Pic is the operator that assigns to X its picard group, usage \Pic(X)
% \Picardstack_{X/B} denotes the Picard stack of X over B
% \Picardfunctor_{X/B} denotes the Picard functor of X over B
\def\Pic{\mathop{\mathrm{Pic}}\nolimits}
\def\Picardstack{\mathcal{P}\!\mathit{ic}}
\def\Picardfunctor{\mathrm{Pic}}
\def\Deformationcategory{\mathcal{D}\!\mathit{ef}}

% --- appended 2026-10-12 (Chief): the project own definitions of the reference/slogan/history
% display environments, copied mechanically from the shared Stacks preamble lines 8-12
% (\usepackage{verbatim} plus the three \newenvironment lines, which the chapter preamble keeps
% verbatim from the project preamble). Nothing added, removed or redefined: these environments
% carry no mathematical content and the project itself typesets them as comments. Without these
% three lines a counted statement whose own text uses them fails with Environment slogan undefined.
\usepackage{verbatim}
\newenvironment{reference}{\comment}{\endcomment}
\newenvironment{slogan}{\comment}{\endcomment}
\newenvironment{history}{\comment}{\endcomment}
\usepackage{cleveref}
\crefname{originallemma}{original lemma}{original lemmas}
\Crefname{originallemma}{Original lemma}{Original lemmas}
\makeatletter
\newlabel{lemma-c-square}{{0F16}{source}{}{sourcepointer-1}{}}
\newlabel{lemma-c-square@cref}{{[originallemma][0][]0F16}{[1][0][]source}{}{}{}}
\makeatother
\crefname{originalsection}{original section}{original sections}
\Crefname{originalsection}{Original section}{Original sections}
\makeatletter
\newlabel{section-unbounded}{{07A5}{source}{}{sourcepointer-2}{}}
\newlabel{section-unbounded@cref}{{[originalsection][0][]07A5}{[1][0][]source}{}{}{}}
\makeatother
\crefname{originalsection}{original section}{original sections}
\Crefname{originalsection}{Original section}{Original sections}
\makeatletter
\newlabel{sites-modules-section-free-abelian-sheaf}{{03CQ}{source}{}{sourcepointer-3}{}}
\newlabel{sites-modules-section-free-abelian-sheaf@cref}{{[originalsection][0][]03CQ}{[1][0][]source}{}{}{}}
\makeatother
\title{A blow-up square that is a pushout with an injective leg produces a distinguished triangle}
\author{The Stacks Project Authors}
\date{stacks-project:sites-cohomology.tex:local-snapshot-20261011; GNU FDL 1.2 or later}
\begin{document}
\maketitle
\noindent\textit{Editorial note.} All delivered mathematical text is the author's own, copied from the pinned Stacks Project chapter file sites-cohomology.tex; no statement or proof step was rewritten, reordered or removed.\par\medskip
\noindent\textit{Editorial note.} The proof is the author's complete argument: the reduction to injective sheaves is stated, and exactness at each of the three positions of the displayed sequence is proved, including the amalgamation step in the middle.\par\medskip
\noindent\textit{Editorial note.} Every labelled reference occurring inside the counted spans is carried in assembly.external\_references as a non-load-bearing navigation pointer (same-chapter ones at their definition site, cross-chapter ones at the invocation inside the counted span), so no load-bearing mathematics is replaced by a stub. No bibliography citation occurs inside the counted spans.\par\medskip
\noindent\textit{Editorial note.} The checked chapter list for same-name collisions is the set of stacks chapters staged in this repository: chow, cohomology, divisors, etale-cohomology, flat, more-algebra, more-morphisms, morphisms, spaces-more-morphisms, varieties together with this card's chapter sites-cohomology.\par\medskip
\noindent\textit{Editorial note.} The preamble is the chapter's own retained minimal preamble (sources/stacks-sites-cohomology/retained-preamble.tex), a verbatim excerpt of the native preamble: the Xy-pic setup and the native theorem declarations with the mathematical macros.\par\medskip
\noindent\textit{Allowed standard prerequisites (editorial).}\begin{itemize}
\item The canonical map of a blow-up square. Same chapter, Lemma lemma-c-square. Supplies the construction whose output is shown to be a distinguished triangle.
\item K-injective representatives in the unbounded derived category. Same chapter, Section section-unbounded. Allows the given complex to be replaced by a K-injective complex of injective sheaves.
\item Free abelian sheaves represented by sheafified Yoneda functors. Modules on Sites, Section sites-modules-section-free-abelian-sheaf. Supplies the description of the values of an injective sheaf as homomorphisms from the sheafified representable functor.
\end{itemize}\medskip
\ifdefined\hypertarget\hypertarget{sourcepointer-1}{}\fi
\noindent\textit{Source reference.} Original lemma 0F16: sources/stacks-sites-cohomology/pinned-original-source.tex, source lines 6001--6001. Navigation pointer invoked at proof line 6080; the counted span contains the pointer only and the cited result is not reproduced inside it, so the pointer is registered as a non-load-bearing navigation stub in the established form for Stacks material. No bibliography citation is involved.\par\medskip
\ifdefined\hypertarget\hypertarget{sourcepointer-2}{}\fi
\noindent\textit{Source reference.} Original section 07A5: sources/stacks-sites-cohomology/pinned-original-source.tex, source lines 3921--3921. Navigation pointer invoked at proof line 6093; the counted span contains the pointer only and the cited result is not reproduced inside it, so the pointer is registered as a non-load-bearing navigation stub in the established form for Stacks material. No bibliography citation is involved.\par\medskip
\ifdefined\hypertarget\hypertarget{sourcepointer-3}{}\fi
\noindent\textit{Source reference.} Original section 03CQ: sources/stacks-sites-cohomology/pinned-original-source.tex, source lines 6112--6112. Navigation pointer invoked at proof line 6112; the counted span contains the pointer only and the cited result is not reproduced inside it, so the pointer is registered as a non-load-bearing navigation stub in the established form for Stacks material. No bibliography citation is involved.\par\medskip
Let $(\mathcal{C}, \mathcal{O})$ be a ringed site. Consider a
commutative diagram
$$
\xymatrix{
E \ar[d] \ar[r] & Y \ar[d] \\
Z \ar[r] & X
}
$$
in the category $\mathcal{C}$. In this situation, given an object
$K$ of $D(\mathcal{O})$ we get what looks like the
beginning of a distinguished triangle
$$
R\Gamma(X, K) \to
R\Gamma(Z, K) \oplus
R\Gamma(Y, K) \to
R\Gamma(E, K)
$$
In the following lemma we make this more precise.

\begin{lemma}
\label{lemma-square-triangle}
In the situation above assume
\begin{enumerate}
\item $h_X^\# = h_Y^\# \amalg_{h_E^\#} h_Z^\#$, and
\item $h_E^\# \to h_Y^\#$ is injective.
\end{enumerate}
Then the construction of Lemma \ref{lemma-c-square}
produces a distinguished triangle
$$
R\Gamma(X, K) \to
R\Gamma(Z, K) \oplus
R\Gamma(Y, K) \to
R\Gamma(E, K) \to R\Gamma(X, K)[1]
$$
functorial for $K$ in $D(\mathcal{C})$.
\end{lemma}

\begin{proof}
We can represent $K$ by a K-injective complex whose terms are injective
abelian sheaves, see Section \ref{section-unbounded}.
Thus it suffices to show: if $\mathcal{I}$ is an injective abelian
sheaf, then
$$
0 \to \mathcal{I}(X) \to
\mathcal{I}(Z) \oplus \mathcal{I}(Y) \to
\mathcal{I}(E) \to 0
$$
is a short exact sequence. The first arrow is injective
because by condition (1) the map $h_Y \amalg h_Z \to h_X$
becomes surjective after sheafification, which means that
$\{Y \to X, Z \to X\}$ can be refined by a covering of $X$.
The last arrow is surjective because $\mathcal{I}(Y) \to \mathcal{I}(E)$
is surjective. Namely, we have
$\mathcal{I}(E) = \Hom(\mathbf{Z}_E^\#, \mathcal{I})$,
$\mathcal{I}(Y) = \Hom(\mathbf{Z}_Y^\#, \mathcal{I})$,
the map $\mathbf{Z}_E^\# \to \mathbf{Z}_Y^\#$ is injective
by (2), and $\mathcal{I}$ is an injective abelian sheaf.
Please compare with Modules on Sites, Section
\ref{sites-modules-section-free-abelian-sheaf}.
Finally, suppose we have $s \in \mathcal{I}(Y)$ and
$t \in \mathcal{F}(Z)$ mapping to the same element of
$\mathcal{I}(E)$. Then $s$ and $t$ define a map
$$
s \amalg t : h_Y^\# \amalg h_Z^\# \longrightarrow \mathcal{I}
$$
which by assumption factors through $h_Y^\# \amalg_{h_E^\#} h_Z^\#$.
Thus by assumption (1) we obtain a unique map
$h_X^\# \to \mathcal{I}$ which corresponds to an element
of $\mathcal{I}(X)$ restricting to $s$ on $Y$ and $t$ on $Z$.
\end{proof}

\end{document}
